\section{Orbital necklace collapse}
\label{sec:orbital}

We now give the strongest favorable calculation after explicitly changing
the primitive unit from full-tree periods to positive-height group conjugacy
classes.  The fixed-height congruence is consistent with current conjugacy
theory for soluble Baumslag--Solitar groups
\citep{ciobanu2020conjugacy,guo2026conjugacy}.

\subsection{Conjugacy at fixed height}

For $r\ge2$, write
\begin{equation}
  \Gr\cong\mathbb Z[1/r]\rtimes\mathbb Z,
  \qquad
  (a,\ell)(b,k)=(a+r^\ell b,\ell+k).
  \label{eq:semidirect}
\end{equation}
The inverse is $(-r^{-\ell}a,-\ell)$, and direct multiplication yields
\begin{equation}
  (c,m)(b,k)(c,m)^{-1}
  =\bigl(r^m b+(1-r^k)c,k\bigr).
  \label{eq:conjugation}
\end{equation}
Height is therefore a conjugacy invariant.  At fixed $k>0$, addition by
$(1-r^k)c$ passes to
\[
  \mathbb Z[1/r]/(r^k-1)\mathbb Z[1/r]
  \cong \mathbb Z/(r^k-1)\mathbb Z,
\]
and conjugation by height $m$ acts by multiplication by $r^m$.

\begin{proposition}[Burnside class count]
\label{prop:burnside}
The number of conjugacy classes at height $k$ is
\begin{equation}
  \Cclass(k)=\frac1k\sum_{j=0}^{k-1}
  \bigl(r^{\gcd(j,k)}-1\bigr)=\Nneck(k)-1,
  \label{eq:class-count}
\end{equation}
where $\Nneck(k)$ is the number of length-$k$ necklaces on $r$ symbols.
\end{proposition}

\begin{proof}
Multiplication by $r$ has exact order $k$ modulo $r^k-1$.  Its $j$th power
fixes
\[
  \gcd(r^j-1,r^k-1)=r^{\gcd(j,k)}-1
\]
residues.  Burnside's lemma gives the first equality.  The ordinary necklace
formula has the same average with $r^{\gcd(j,k)}$ in place of
$r^{\gcd(j,k)}-1$, proving the second equality.
\end{proof}

The explicit digit map sends $(d_0,\ldots,d_{k-1})$ to
$\sum_i d_ir^i$ modulo $r^k-1$.  Rotation acts by multiplication by $r$.
The only duplicate in the integer range $[0,r^k-1]$ is the pair
$0^k$ and $(r-1)^k$, which explains the subtraction of one.

\begin{figure}[t]
  \centering
  \resizebox{0.99\textwidth}{!}{\begin{tikzpicture}[
  >=Latex,
  node distance=6mm,
  box/.style={rounded corners,draw=charcoal,thick,align=center,
    fill=white,inner sep=5pt,font=\small,text width=3.25cm,
    minimum height=1.15cm},
  algebra/.style={box,draw=formalblue},
  count/.style={box,draw=deepgreen},
  result/.style={box,draw=warningamber},
  arr/.style={->,very thick,formalblue}
]
\node[algebra] (semi)
  {$\mathbb Z[1/r]\rtimes\mathbb Z$\\
   $(c,m)(b,k)(c,m)^{-1}$\\
   $=(r^m b+(1-r^k)c,k)$};
\node[algebra,right=of semi] (residue)
  {fixed height $k>0$\\
   $\mathbb Z/(r^k-1)\mathbb Z$\\
   modulo $x\mapsto rx$};
\node[count,right=of residue] (necklace)
  {$r$-ary necklaces\\cyclic rotation\\
   $0^k\sim(r-1)^k$};
\node[count,right=of necklace] (burnside)
  {$C_r(k)=N_r(k)-1$\\
   primitive $P_r(k)$\\
   by M\"obius inversion};
\node[result,right=of burnside] (zeta)
  {$Z_{+,r}(z)=\dfrac{1-z}{1-rz}$\\
   $z\mapsto r^{-s}z$\\
   under $\Delta^{-s}$};

\draw[arr] (semi) -- (residue);
\draw[arr] (residue) -- (necklace);
\draw[arr] (necklace) -- (burnside);
\draw[arr] (burnside) -- (zeta);

\node[rounded corners,draw=stopred,dashed,thick,fill=white,
  below=9mm of necklace,align=center,font=\small,text=stopred,
  text width=10.2cm] (barrier)
  {exact group-orbital classification, but not a periodic point or an
   ordinary Fredholm factor of the literal full-tree shift};
\draw[dashed,very thick,stopred,-|] (barrier.north) -- (necklace.south);
\end{tikzpicture}
}
  \caption{The exact positive-height group-orbital pipeline.  Conjugacy
  reduces to multiplication-by-$r$ residue orbits, which are endpoint-
  identified necklaces.  Burnside and primitive-root inversion give the
  rational product.  The dashed barrier records the decisive ownership fact:
  this group-conjugacy ledger is not a periodic ledger or ordinary Fredholm
  factor of the literal full-tree shift.}
  \label{fig:orbital-necklace}
\end{figure}

\subsection{Primitive roots and the Euler product}

Let $\Pclass(k)$ count primitive positive-height conjugacy classes.  A word of
minimal period $d\mid k$ repeats $k/d$ times, and the associated translation
coordinate is multiplied by
$1+r^d+\cdots+r^{k-d}$.  This is exactly the first coordinate of the
corresponding semidirect-product power.  M\"obius inversion gives
\begin{equation}
  \Pclass(1)=r-1,
  \qquad
  \Pclass(k)=\frac1k\sum_{d\mid k}\mu(d)r^{k/d}\quad(k>1),
  \qquad
  \Cclass(k)=\sum_{d\mid k}\Pclass(d).
  \label{eq:primitive-count}
\end{equation}

The ordinary Witt necklace product is $(1-rz)^{-1}$.  Our primitive counts
differ only by one missing degree-one class, so
\begin{equation}
  \Zpos(z)
  =\prod_{k\ge1}(1-z^k)^{-\Pclass(k)}
  =\frac{1-z}{1-rz}.
  \label{eq:orbital-product}
\end{equation}
This factorization owns primitive/repetition bookkeeping inside the
group-orbital ledger.  It does not repair the object change.

\subsection{The modular cocycle only rescales length}

At positive height $k$, Convention~\ref{conv:height} gives
$\modc^{-s}=r^{-sk}$.  Weighting \eqref{eq:orbital-product} therefore makes
the substitution $z\mapsto r^{-s}z$:
\begin{equation}
  Z_{+,r,s}(z)
  =\frac{1-r^{-s}z}{1-r^{1-s}z}.
  \label{eq:modular-product}
\end{equation}
The cocycle supplies no second label beyond signed tree length.  In
particular, it produces neither a graded cancellation nor a prime/composite
split.

\subsection{Balanced divergence and genericity}

For $r=1$, the group is $\mathbb Z^2$ and conjugacy is equality.  At each
height $k>0$, the elements $(b,k)$ for $b\in\mathbb Z$ give infinitely many
classes.  The orbital Euler product is not locally finite.  Quotienting the
infinite height-zero kernel leaves a generic line action and changes the
object again.

\begin{table}[t]
\centering
\small
\caption{First six total and primitive orbital counts.  Prime and composite
rows follow one index law; the table is a genericity control, not a
factorization classifier.}
\label{tab:necklace-counts}
\begin{tabular}{cll}
\toprule
$r$ & $C_r(1{:}6)$ & $P_r(1{:}6)$ \\
\midrule
2 & $1,2,3,5,7,13$ & $1,1,2,3,6,9$ \\
3 & $2,5,10,23,50,129$ & $2,3,8,18,48,116$ \\
4 & $3,9,23,69,207,699$ & $3,6,20,60,204,670$ \\
5 & $4,14,44,164,628,2634$ & $4,10,40,150,624,2580$ \\
6 & $5,20,75,335,1559,7825$ & $5,15,70,315,1554,7735$ \\
\bottomrule
\end{tabular}
\end{table}

Every cyclic index-$r$ ascending HNN control reproduces
\eqref{eq:orbital-product}.  The formula remembers the numerical incidence
index, but it does not distinguish the source presentation from its matched
controls.  Under the preregistered rule, that genericity is terminal.
